\chapter{Marco Teórico}

Este capítulo establece los fundamentos conceptuales, neurofisiológicos y matemáticos que sustentan el desarrollo de esta tesis. En una primera instancia, se describe la naturaleza de la señal electroencefalográfica (EEG) y las dificultades inherentes a su captura. Posteriormente, se formaliza la teoría del Procesamiento de Señales en Grafos (GSP) como marco de resolución, detallando las arquitecturas de interpolación clásica frente a las estrategias espacio-temporales modernas. Finalmente, se definen rigurosamente las métricas utilizadas para evaluar la fidelidad de la reconstrucción.

\section{Fundamentos Neurofisiológicos de la Electroencefalografía (EEG)}

El electroencefalograma es el registro de la actividad eléctrica del cerebro captada desde la superficie del cuero cabelludo. Esta señal refleja la suma espacio-temporal de los potenciales postsinápticos excitatorios e inhibitorios generados por decenas de miles de neuronas piramidales sincronizadas en la corteza cerebral. Debido a que las fuentes de estas corrientes eléctricas se encuentran separadas de los electrodos por diversas capas de tejido (piamadre, líquido cefalorraquídeo, cráneo y cuero cabelludo), la señal sufre un efecto de filtrado espacial y atenuación conocido como \textit{efecto de conducción de volumen} (volume conduction).

En paralelo al uso clínico tradicional, la literatura reciente en procesamiento de EEG ha enfatizado que los retos actuales no se reducen al filtrado de ruido, sino a la preservación simultánea de estructura espacial, dinámica temporal y contenido espectral en presencia de datos incompletos \cite{sharma_emerging_2024}. Esta necesidad ha impulsado la transición desde modelos puramente geométricos hacia marcos de señal sobre grafos con regularización física y estadística \cite{calazans_machine_2024, li_graph_2023}.

Este efecto provoca que el potencial eléctrico generado por una única fuente cortical se propague e irradie, siendo detectado por múltiples electrodos simultáneamente. Como resultado, las señales registradas en electrodos adyacentes presentan un alto grado de correlación espacial. Sin embargo, esta propagación no es perfectamente esférica ni homogénea debido a las anisotropías en la conductividad del cráneo. Por ello, si bien la distancia euclidiana entre electrodos es informativa, no es una representación perfecta de la conectividad funcional subyacente.

Las señales EEG se caracterizan por su naturaleza oscilatoria y transitoria. Clásicamente, la actividad continua se descompone en bandas de frecuencia (Delta, Theta, Alfa, Beta, Gamma), mientras que la actividad evocada o fásica se manifiesta como Potenciales Relacionados a Eventos (ERP, por sus siglas en inglés). Los ERPs, como el N100 o el P300, son fluctuaciones de voltaje transitorias y de fase bloqueada que reflejan el procesamiento cognitivo o sensorial. La preservación de la morfología y la latencia temporal de estos ERPs es crítica en aplicaciones clínicas y de interfaces cerebro-computador (BCI).

\section{El Problema de Canales Faltantes e Interpolación Clásica}

En la práctica, la adquisición de datos EEG es altamente susceptible a la pérdida parcial de información. Los canales pueden resultar inutilizables por mala conductividad del gel, desprendimiento físico del electrodo, o saturación del amplificador por artefactos de movimiento. Matemáticamente, si la actividad cerebral medida se representa como un vector $\mathbf{x} \in \mathbb{R}^N$ para $N$ electrodos, un evento de desconexión implica que solo observamos un subconjunto de electrodos $\mathcal{K} \subset \{1, \dots, N\}$. El objetivo de la interpolación es estimar los valores en el conjunto de canales faltantes $\mathcal{U}$, tal que $\mathcal{K} \cup \mathcal{U} = \mathcal{V}$ y $\mathcal{K} \cap \mathcal{U} = \emptyset$.

\subsection{Interpolación Geométrica y Estadística}

Clásicamente, el problema ha sido resuelto proyectando los valores conocidos mediante bases espaciales dependientes de la distancia euclidiana. El método estándar \textit{de facto} es la interpolación por \textbf{Splines Esféricos} (\textit{Spherical Splines}) \cite{perrin_spherical_1989}. Este método asume la topología de la cabeza como una esfera unitaria y proyecta el potencial en una ubicación faltante $\mathbf{r}$ como:
\begin{equation}
    V(\mathbf{r}) = c_0 + \sum_{i \in \mathcal{K}} c_i \, g(\mathbf{r}, \mathbf{r}_i),
\end{equation}
donde los coeficientes $c_i$ se obtienen resolviendo un sistema lineal sobre los canales conocidos, y $g(\cdot, \cdot)$ es una función de base que decae con la distancia angular entre los vectores posición de los electrodos. Análogamente, las Funciones de Base Radial (RBF) \cite{jager_reconstruction_2016} o la interpolación por vecinos más cercanos asumen suavidad espacial paramétrica en 3D.

\subsection{Implementación de referencia en MNE-Python: \code{interpolate\_bads}}

En la práctica neurocientífica moderna, una referencia operacional importante no es solo la formulación matemática original de Perrin et al., sino su implementación en librerías ampliamente usadas. En MNE-Python, \code{Raw.interpolate\_bads()} opera sobre canales marcados como malos y, para canales EEG, usa por defecto \code{method='spline'} cuando no se especifica otro método \cite{gramfort2013meg, gramfort_mne_software_2014, mne_interpolate_bads_docs}. Es importante distinguir esta opción de \code{method='MNE'}: en la API de MNE, \code{'MNE'} alude a una proyección minimum-norm sobre una esfera, mientras que la interpolación EEG por defecto corresponde a splines esféricos.

La implementación \code{spline} de MNE mantiene el núcleo conceptual de Perrin et al.~\cite{perrin_spherical_1989}: proyecta las posiciones de los electrodos sobre una esfera, calcula cosenos angulares entre sensores y construye la matriz $G$ mediante una expansión en polinomios de Legendre. No obstante, introduce decisiones implementacionales que la vuelven un baseline particularmente fuerte y estable en la práctica:
\begin{enumerate}
  \item \textbf{Normalización geométrica y origen esférico}: las posiciones se trasladan respecto de un origen de esfera, por defecto ajustado automáticamente a la digitización de cabeza cuando está disponible, y luego se normalizan a norma unitaria antes de calcular los ángulos.
  \item \textbf{Expansión truncada y rigidez fija}: MNE evalúa la función $g$ con rigidez $m=4$ y 50 términos de Legendre, una aproximación finita eficiente del desarrollo infinito usado en la teoría original.
  \item \textbf{Regularización diagonal}: antes de invertir el sistema, añade un término pequeño $\alpha=10^{-5}$ a la diagonal de $G$. Este término mejora el condicionamiento numérico cuando hay electrodos cercanos, montajes densos o configuraciones casi degeneradas.
  \item \textbf{Pseudoinversa y término constante}: resuelve el sistema aumentado con una pseudoinversa de Moore--Penrose e incluye el término constante de la formulación spline, reduciendo fallos por singularidad exacta.
  \item \textbf{Controles de calidad}: usa solo canales buenos como base de interpolación, advierte si el ajuste esférico es pobre y permite manejar posiciones inválidas con error, warning o valores \code{NaN}.
\end{enumerate}

Estas decisiones no constituyen un método nuevo separado de las splines esféricas; son una adaptación computacional robusta de la formulación clásica a datos reales, montajes heterogéneos y flujos de trabajo reproducibles. Por ello, \code{MNE interpolate\_bads(method='spline')} debe tratarse como un baseline clásico fuerte: es rápido, tiene pocos grados de libertad, explota todos los canales buenos disponibles y está muy alineado con la hipótesis física de conducción de volumen, que suaviza espacialmente los potenciales de cuero cabelludo.

El problema central de las familias geométricas es que calculan la interpolación operando sobre cada instante de tiempo $t$ de forma aislada. Al carecer de \textit{memoria} temporal, son incapaces de preservar la fase de transitorios rápidos o la dinámica espectral subyacente, introduciendo discontinuidades y ruido artificial de alta frecuencia en el canal interpolado.

Esta limitación ha sido señalada también en estudios de interpolación no-suave sobre grafos, donde se evidencia que la hipótesis de suavidad global puede ser insuficiente para señales con transitorios locales marcados \cite{mazarguil_non-smooth_2022}. En el contexto específico de EEG, reportes aplicados han mostrado que incorporar estructura de grafo mejora la imputación frente a enfoques clásicos, pero con fuerte dependencia de la calidad del grafo y su parametrización \cite{weinstein_imputing_2023, jiang_graph-based_2021}.

\section{Teoría de Grafos y Procesamiento de Señales en Grafos (GSP)}

Para abandonar la restricción de la esfera euclidiana, el Procesamiento de Señales en Grafos (GSP) \cite{shuman_emerging_2013, ortega_graph_2018-1, ortega_introduction_2022, leus_graph_2023} propone modelar los electrodos y sus interacciones complejas como un grafo ponderado. 

Desde una perspectiva histórica, esta transición no solo introduce nuevas herramientas matemáticas, sino también un cambio epistemológico: la señal ya no se interpreta como una función en una superficie continua idealizada, sino como una realización en una red irregular donde topología y señal se co-determinan. Esta visión ha demostrado utilidad en aplicaciones biomédicas para representación espectral, reducción de dimensionalidad y análisis de conectividad funcional \cite{rui_dimensionality_2016, miri_spectral_2024}.

Un grafo no dirigido se define como $\mathcal{G} = (\mathcal{V}, \mathcal{E}, \mathbf{W})$, donde $\mathcal{V}$ es el conjunto de $N$ vértices (electrodos), $\mathcal{E}$ es el conjunto de aristas y $\mathbf{W} \in \mathbb{R}^{N \times N}$ es la matriz de adyacencia ponderada. El peso $W_{ij} \ge 0$ codifica la fuerza de la relación (física o funcional) entre el electrodo $i$ y el electrodo $j$. Una señal de EEG en un instante determinado se define como una ``señal sobre el grafo'' $\mathbf{x} \in \mathbb{R}^N$, una función que mapea cada vértice a un escalar (voltaje).

\subsection{El Laplaciano del Grafo}

El operador fundamental en GSP es el Laplaciano Combinatorial $\mathbf{L} = \mathbf{D} - \mathbf{W}$, donde $\mathbf{D}$ es la matriz diagonal de grados con $D_{ii} = \sum_j W_{ij}$. Al ser simétrica y semidefinida positiva, $\mathbf{L}$ posee un conjunto completo de autovectores reales ortonormales $\mathbf{U} = [\mathbf{u}_1, \dots, \mathbf{u}_N]$ y autovalores no negativos $\boldsymbol{\Lambda} = \text{diag}(\lambda_1, \dots, \lambda_N)$, tal que $\mathbf{L} = \mathbf{U} \boldsymbol{\Lambda} \mathbf{U}^\top$ con $0 = \lambda_1 \le \lambda_2 \dots \le \lambda_N$.

La forma cuadrática del Laplaciano cuantifica la ``variación total'' o la suavidad espacial de la señal $\mathbf{x}$ sobre la topología de la red:
\begin{equation}
    S(\mathbf{x}) = \mathbf{x}^\top \mathbf{L} \mathbf{x} = \frac{1}{2} \sum_{(i,j) \in \mathcal{E}} W_{ij} (x_i - x_j)^2.
    \label{eq:smoothness}
\end{equation}
Un valor bajo de $S(\mathbf{x})$ indica que los valores de señal entre electrodos fuertemente conectados son muy similares, lo que representa una señal suave o de baja frecuencia en el dominio espacial.

En términos de marco conceptual general, la formulación moderna de señales sobre grafos y sus operadores espectrales se apoya también en revisiones fundacionales complementarias ampliamente citadas en la literatura reciente \cite{ortega_graph_2018}.

\subsection{Transformada de Fourier sobre Grafos (GFT)}

Por analogía con el análisis de Fourier clásico, los autovectores de $\mathbf{L}$ conforman la base de Fourier del grafo, y los autovalores $\lambda_l$ actúan como las ``frecuencias''. La Transformada de Fourier sobre Grafos (GFT) de una señal $\mathbf{x}$ se define como su proyección sobre los autovectores:
\begin{equation}
    \hat{\mathbf{x}} = \mathbf{U}^\top \mathbf{x}.
\end{equation}
En el EEG, los componentes asociados a autovalores pequeños ($\lambda \approx 0$) representan patrones de actividad topográfica generalizada y sincrónica (ej., ritmo alfa global). Los autovalores grandes representan variaciones espaciales abruptas (ej., ruido electromagnético de alta frecuencia o actividad muscular aislada focal). Así, la interpolación espacial basada en GSP actúa intrínsecamente como un filtro paso-bajo direccional sobre la superficie del cráneo \cite{goerttler_understanding_2023}.

\subsection{Construcción y Aprendizaje de Grafos}

El éxito del GSP recae enteramente en la matriz $\mathbf{W}$. En neurofisiología, la topología puede diseñarse utilizando tres estrategias generales. La primera consiste en construir \textbf{grafos geométricos} \textit{a priori}, usando la inversa de la distancia euclidiana entre las coordenadas 3D de los electrodos. Una variante popular es el Grafo $k$-Vecinos Más Cercanos Gaussianos ($k$-NN Gaussiano), donde $W_{ij} = \exp(-d(i,j)^2/\sigma^2)$ si $j$ está entre los $k$ vecinos de $i$, y $0$ en caso contrario; Tamaru et al.~\cite{tamaru_optimizing_2024} han explorado variantes de $k$ dinámico para combatir irregularidades locales en la densidad de mallas. La segunda estrategia emplea \textbf{grafos funcionales o de correlación}, extraídos empíricamente de la covarianza de la señal temporal o mediante modelos adaptativos como \textit{Adaptive Edge Weighting}~\cite{karasuyama_adaptive_2017}. Finalmente, una tercera familia corresponde al \textbf{aprendizaje de grafos} (\textit{Graph Learning}), donde se plantean problemas de optimización inversos que buscan inferir la matriz $\mathbf{W}$ directamente desde los datos, exigiendo simultáneamente que la red resultante sea esparsa y que las señales observadas exhiban alta suavidad laplaciana~\cite{kalofolias_how_2016, qiao_data-driven_2018}.

Dentro de esta línea, los trabajos sobre aprendizaje de grafos de gran escala y selección de vecindad para señales sobre grafos han reforzado criterios prácticos de construcción topológica en escenarios con alta dimensionalidad y heterogeneidad de conectividad \cite{kalofolias_large_2019, shekkizhar_graph_2020, shekkizhar_neighborhood_2023}.

\section{Estrategias de Reconstrucción de Grafos}

El problema de interpolación se formula como un problema inverso mal propuesto. Definimos una matriz diagonal de muestreo $\mathbf{S} \in \mathbb{R}^{N \times N}$ donde $S_{ii} = 1$ si $i \in \mathcal{K}$ (observado) y $0$ si $i \in \mathcal{U}$ (faltante). Observamos la señal parcial $\mathbf{y} = \mathbf{S} \mathbf{x}$.

\subsection{Regularización Tikhonov Espacial}
El enfoque clásico de GSP es la regularización de Tikhonov sobre grafos \cite{narang_signal_2013}. Asumiendo ruido gaussiano, se busca recuperar $\mathbf{x}$ minimizando el error de los nodos conocidos más una penalización a la variación total (Ecuación \ref{eq:smoothness}):
\begin{equation}
    \hat{\mathbf{x}} = \arg\min_{\mathbf{z}} \| \mathbf{y} - \mathbf{S} \mathbf{z} \|_2^2 + \alpha (\mathbf{z}^\top \mathbf{L} \mathbf{z}).
    \label{eq:tikhonov_general}
\end{equation}
donde el hiperparámetro $\alpha > 0$ controla el nivel de suavidad forzada (\textit{spatial smoothness}). 

Además de la formulación cerrada, existen variantes iterativas localizadas que reducen costo computacional y facilitan despliegue en redes de sensores de mayor escala \cite{narang_localized_2013}.

\subsection{Recuperación Bandlimited y Espacios RKHS (BGSRP)}
Un enfoque alternativo, BGSRP (\textit{Bandlimited Graph Signal Recovery in RKHS}) \cite{zhang_recovery_2024}, asume que las señales EEG son estrictamente limitadas en banda (solo contienen frecuencias de grafo $\lambda \le \lambda_c$). Utilizando un Espacio de Hilbert de Núcleo Reproductor (RKHS), se reconstruye la señal proyectando sobre una matriz núcleo $\mathbf{K}$ ensamblada a partir de los primeros $B$ autovectores. En paralelo, la literatura reciente también ha reforzado la base teórica de la recuperación sobre grafos con resultados de garantía de identificación y acotación de coherencia \cite{morgenstern_theoretical_2025}, además de extender los marcos de reconstrucción hacia formulaciones generalizadas vértice-tiempo \cite{zhao_generalized_2025}.

En la misma línea, los desarrollos más recientes sobre grafos variables en el tiempo muestran que, cuando la estructura de dependencia evoluciona junto con la señal, los modelos con acoplamiento espacio-temporal explícito ofrecen mayor capacidad de generalización para tareas de reconstrucción y predicción \cite{yan_signal_2026, krishnan_learning_2026}. Esta evidencia respalda directamente el uso de regularización temporal en escenarios EEG con pérdida severa o contigua de canales.

\subsection{Regularización Espacio-Temporal (TRSS)}
Todos los métodos anteriores reconstruyen una muestra temporal a la vez. Sin embargo, un EEG es una matriz continua $\mathbf{X} \in \mathbb{R}^{N \times T}$. La Regularización de Suavidad de Sobolev Temporal (TRSS, \textit{Temporal Regularized Sobolev Smoothness}) \cite{giraldo2022} revoluciona la interpolación minimizando simultáneamente los gradientes espaciales y las derivadas temporales:
\begin{equation}
    \hat{\mathbf{X}} = \arg\min_{\mathbf{Z}} \| \mathbf{S}\mathbf{X} - \mathbf{S}\mathbf{Z} \|_F^2 + \alpha \, \text{tr}(\mathbf{Z}^\top \mathbf{L} \mathbf{Z}) + \beta \| \mathbf{Z}\mathbf{D}_t \|_F^2,
    \label{eq:trss_obj}
\end{equation}
donde $\mathbf{D}_t$ es una matriz que actúa como operador de diferencias finitas temporales (i.e. penalizando cambios bruscos en el tiempo). Este \textit{prior} temporal, controlado por $\beta$, es excepcionalmente poderoso para el EEG, ya que refleja matemáticamente la inercia neurofisiológica de los generadores dipolares del cerebro.

Para evaluar rigurosamente la calidad de la reconstrucción se emplean métricas de error escalar, morfológico y espectral. Sean $\mathbf{x}$ el vector real de un canal en el tiempo de longitud $T$ y $\hat{\mathbf{x}}$ su reconstrucción.

En el dominio de la amplitud, el \textbf{Mean Absolute Error} (MAE) mide la desviación absoluta media, resultando altamente intuitivo para bioseñales:
\begin{equation}
    \text{MAE}(\mathbf{x}, \hat{\mathbf{x}}) = \frac{1}{T} \sum_{t=1}^T | x(t) - \hat{x}(t) |
\end{equation}
El \textbf{Root Mean Square Error} (RMSE) añade una penalización severa a las grandes desviaciones aisladas, sirviendo como indicador de inestabilidades numéricas en la interpolación:
\begin{equation}
    \text{RMSE}(\mathbf{x}, \hat{\mathbf{x}}) = \sqrt{\frac{1}{T} \sum_{t=1}^T (x(t) - \hat{x}(t))^2}
\end{equation}
Por su parte, el \textbf{Signal-to-Noise Ratio} (SNR) cuantifica cuánta energía de la señal original es preservada frente a la energía del error residual, expresado en decibelios:
\begin{equation}
    \text{SNR}(\mathbf{x}, \hat{\mathbf{x}}) = 10 \log_{10} \left( \frac{\sum x(t)^2}{\sum (x(t) - \hat{x}(t))^2} \right)
\end{equation}

En el dominio temporal-morfológico, el \textbf{Dynamic Time Warping} (DTW) busca la alineación temporal óptima de secuencias que minimice la distancia euclidiana. A diferencia del MAE o RMSE, que comparan punto a punto $\hat{x}(t)$ con $x(t)$, DTW permite medir la fidelidad morfológica (la ``forma'' del ERP) sin ser excesivamente penalizado por minúsculos desfases de latencia de fase, lo que lo convierte en una herramienta especialmente valiosa para el análisis de EEG.

En el dominio espectral, la \textbf{Log Spectral Distance} (LSD) mide la distancia entre las Densidades Espectrales de Potencia de ambas señales en escala logarítmica, asegurando que la reconstrucción preserve la distribución de energía en las bandas Alfa, Beta y Gamma, aspecto crucial para aplicaciones de BCI. Complementariamente, la \textbf{Magnitude-Squared Coherence} (Coherencia) proporciona una medida entre $0$ y $1$ que indica el nivel de correlación de fase lineal entre ambas señales a través de distintas frecuencias.

\section{Brechas en el Estado del Arte}

Si bien estudios recientes han aplicado GSP para tareas avanzadas como extracción de características para desórdenes de conciencia \cite{mortaheb_graph_2019} o detección de Alzheimer \cite{mootoo_detecting_2023}, la aplicación de grafos dinámicos en la recuperación de señales \cite{yan_signal_2026, krishnan_learning_2026} ha puesto en evidencia una grave falla metodológica en la evaluación biomédica. Específicamente, en el problema de canales EEG, los \textit{baselines} como Splines Esféricos o RBF suelen ejecutarse con hiperparámetros estáticos o predeterminados por el software original, en tanto que los métodos nuevos sí se calibran. Esto enmascara el rendimiento asintótico real y entrega una ventaja desleal al nuevo algoritmo \cite{calazans_machine_2024}. Esta tesis viene a sanear dicho vacío, sometiendo tanto el estado del arte espacial como los modelos temporales TRSS a un proceso de sintonización bayesiana irrestricto, con un marco métrico fisiológico inusualmente riguroso y un fuerte énfasis en la mitigación de fallas contiguas (\textit{nearby masking}).

\section{Criterios para una comparación justa en reconstrucción EEG}
Una comparación metodológica en EEG es justa solo si controla al menos cuatro fuentes de variabilidad: la densidad espacial de sensores, el modo de pérdida de canales, el preprocesamiento temporal y la sintonización de hiperparámetros. La densidad espacial determina cuánta información vecinal queda disponible cuando se elimina un subconjunto de electrodos. El modo de pérdida distingue fallas independientes de fallas contiguas asociadas a desconexión de una región de la gorra. El preprocesamiento temporal fija el ancho de banda de la señal sobre la cual se juzga la reconstrucción. Finalmente, la sintonización de hiperparámetros evita que un método parezca inferior solo por haber sido evaluado con valores por defecto.

Este criterio es especialmente importante al comparar métodos de distinta naturaleza. Un spline es rápido y geométricamente interpretable; un método GSP instantáneo introduce una topología discreta; un método temporal como TRSS añade una hipótesis dinámica adicional. La pregunta de investigación no es cuál método gana bajo una configuración arbitraria, sino cuál familia mantiene mejor fidelidad cuando cada una recibe una oportunidad de optimización comparable.
